%% LyX 2.4.4 created this file.  For more info, see https://www.lyx.org/.
%% Do not edit unless you really know what you are doing.
\documentclass[english]{article}
\usepackage[T1]{fontenc}
\usepackage[latin9]{inputenc}
\usepackage{units}
\usepackage{enumitem}
\usepackage{amsmath}
\usepackage{amsthm}
\usepackage{geometry}
\geometry{verbose,tmargin=1in,bmargin=1in,lmargin=1in,rmargin=1in}

\makeatletter
%%%%%%%%%%%%%%%%%%%%%%%%%%%%%% Textclass specific LaTeX commands.
\numberwithin{equation}{section}
\numberwithin{figure}{section}
\newlength{\lyxlabelwidth}      % auxiliary length 

%%%%%%%%%%%%%%%%%%%%%%%%%%%%%% User specified LaTeX commands.
\usepackage{fancyhdr}
\usepackage{lastpage}
\pagestyle{fancy}
\usepackage{enumitem}
\usepackage{ifthen}
\everymath=\expandafter{\the\everymath\displaystyle}
%\DeclareMathSizes{10}{10}{10}{10}
\usepackage{multicol}

\usepackage{tikz}
\usetikzlibrary{patterns}
\usetikzlibrary{plotmarks}
\usepackage{pgfplots}
\pgfplotsset{compat=newest}

\usepackage{calc}
\usepackage[nomessages]{fp}

\usepackage{datetime}
% Added by lyx2lyx
\setlength{\parskip}{\smallskipamount}
\setlength{\parindent}{0pt}

\makeatother

\usepackage{babel}
\begin{document}
\renewcommand{\author}{} %%your name%% Dr.\,James Ashe

\newcommand{\classNum}{232}

\newcommand{\classSec}{} %%class section%%

\newcommand{\examType}{Final Exam}

\renewcommand{\date}{} %\formatdate{15}{5}{2013} %%copy and paste for date {day}{month}{year}

%%First page's header%%%%%%%%%%%%%%%%%%%%%%
\fancypagestyle{firststyle} %%header settings for first pagr
{
	\fancyhead[L]{MTH \classNum{} \classSec{} \author{}\\
	\textbf{\examType{}} ~\date{}}
	\fancyhead[C]{}
	\fancyhead[R]{\textbf{Name:\hspace{7cm}}}
	\setlength{\headsep}{20pt}
}
%%%%%%%%%%%%%%%%%%%%%%%%%%%%%%%%%%%%%%%%%%

%%Next page's header%%%%%%%%%%%%%%%%%%%%%%%
\setlist{nolistsep}
\lhead{MTH \classNum{} \classSec{}} 
\chead{} 
\rhead{\textbf{Name:\hspace{7cm}}}
\cfoot{\thepage{}~of~\pageref{LastPage}}
\setlength{\headsep}{10pt} 
%%%%%%%%%%%%%%%%%%%%%%%%%%%%%%%%%%%%%%%%%%%

%%Various settings; see the preamble for other settings%%%%%
%\setenumerate[0]{leftmargin=12pt,itemindent=0pt,labelwidth=0pt}
\setlist{topsep=.5pt}
\renewcommand{\labelenumi}{\textbf{\arabic{enumi}.}}
\renewcommand{\labelenumii}{\textbf{\Alph{enumii}.}}
\thispagestyle{firststyle}
%%%%%%%%%%%%%%%%%%%%%%%%%%%%%%%%%%%%%%%%%%

\newcommand{\xmax}{0}
\newcommand{\xmin}{-\xmax}
\newcommand{\xstep}{0}
\newcommand{\ymax}{0}
\newcommand{\ymin}{-\ymax}
\newcommand{\ystep}{0} 
\newcommand{\Fx}{x}
\newcommand{\Gx}{x}

\newcommand{\dspace}{\vspace{0cm}}
\newcommand{\samples}{500}
\begin{enumerate}
\item The technique of integration used to find the indefinite integral
${\displaystyle \int x^{2}e^{-x}\,dx}$ is 

\begin{multicols}{4}
\begin{enumerate}
\item partial fractions\textbf{\columnbreak}
\item by parts \textbf{\columnbreak}
\item trigonometric substitution\textbf{\columnbreak}
\item power rule
\end{enumerate}
\end{multicols}\vfill\dspace
\item If $x>1$, and $y={\displaystyle \int_{1}^{x^{2}}\sqrt{1+\sqrt{t}}\,dt}$,
then $\frac{dy}{dx}$ is equal to

\begin{multicols}{4}
\begin{enumerate}
\item \textbf{$1+x$\columnbreak}
\item \textbf{$2x\sqrt{1+x}$\columnbreak}
\item \textbf{$\sqrt{1+\sqrt{x}}$\columnbreak}
\item \textbf{$\sqrt{1+x^{2}}$}
\end{enumerate}
\end{multicols}\vfill\dspace
\item The \emph{length of the arc }on the graph of $y=f(x)$ from the point
$\left(a,f(a)\right)$ to the point $\left(b,f(b)\right)$ is computed
using

\begin{multicols}{4}
\begin{enumerate}
\item ${\displaystyle \int_{a}^{b}\sqrt{1+\left[f'(x)\right]^{2}}}$\columnbreak
\item ${\displaystyle \int_{a}^{b}\sqrt{1+\left[f(x)\right]^{2}}}$\columnbreak
\item ${\displaystyle \pi\int_{a}^{b}\sqrt{1+\frac{d^{2}y}{dx^{2}}}}$\columnbreak
\item ${\displaystyle \pi\int_{a}^{b}\left(f(x)\right)^{2}x}$\columnbreak
\end{enumerate}
\end{multicols}\vfill\dspace
\item The integral expression which gives the area of the shaded region
$R$ in the figure is

\begin{multicols}{3}
\begin{enumerate}
\item ${\displaystyle \int_{a}^{b}\left(g(x)-f(x)\right)\,dx}$\vfill
\item $\pi{\displaystyle \int_{a}^{b}\left(f(x)-g(x)\right)\,dx}$\vfill\columnbreak
\item ${\displaystyle \int_{a}^{b}\left(f(x)-g(x)\right)^{2}\,dx}$\vfill
\item ${\displaystyle \int_{a}^{b}\left(f(x)-g(x)\right)\,dx}$\vfill\columnbreak
\item []\renewcommand{\xmax}{15}
\renewcommand{\xmin}{-4}
\renewcommand{\xstep}{0} %must be xmin+n
\renewcommand{\ymax}{15}
\renewcommand{\ymin}{-3}
\renewcommand{\ystep}{-} %must be ymin+n
%%%%%%%
\begin{tikzpicture}
\begin{axis}[height=4cm, 
	width=5cm, 
	axis lines=middle,
	%axis line style={-latex}, 
	xmin=\xmin,xmax=\xmax, 
	ymin=\ymin,ymax=\ymax,
	%restrict y to domain=-1:10,
	no markers, 
	ticks=none,
	xlabel={$x$},ylabel={$y$}, 
	%every axis x label/.append style={anchor=north}, 
	%every axis y label/.append style={anchor=east}, 
	area style,tension=0.08]

\addplot+[domain=3:12,fill,black!30!white]{1*cos(deg(x/2))+10} \closedcycle;
	%\addplot+[domain=3:12,fill,black!30!white]{1*cos(deg(x/2.7))+3} \closedcycle;
	\addplot+[domain=3:12,fill,white]{1*cos(deg(x/2.7))+3} \closedcycle; 

	\draw[thick,black,dashed] (axis cs:3,0) node[below] {$a$} -- (axis cs:3,10.0707) ;
	\draw[thick,black,dashed] (axis cs:12,0) node[below] {$b$} -- (axis cs:12,10.9602);
	
	\addplot[domain=\xmin+1:\xmax,thick,smooth,red,samples=\samples]{1*cos(deg(x/2))+10} node[above,pos=.04] {$f(x)$}; 
	\addplot[domain=\xmin+1:\xmax,thick,smooth,blue,samples=\samples]{1*cos(deg(x/2.7))+3} node[above,pos=.04] {$g(x)$};

	\node at (axis cs:7.5,6) {$R$};
\end{axis}
\end{tikzpicture}
\end{enumerate}
\end{multicols}\vfill
\item Let $f(x)$ be continuous on $[-5,5]$ and ${\displaystyle \int_{0}^{5}f(x)\,dx}=13$.
If for each $x$ in $[-5,5]$, $f(-x)=f(x)$, i.e., $f(x)$ is an
even function, then ${\displaystyle \int_{-5}^{5}f(x)\,dx}=$

\begin{multicols}{4}
\begin{enumerate}
\item 26\columnbreak
\item 13\columnbreak
\item 0\columnbreak
\item -13\columnbreak
\end{enumerate}
\end{multicols}\vfill\dspace
\item Let $f(x)$ be continuous on $[-5,5]$ and ${\displaystyle \int_{-5}^{0}f(x)\,dx}=13$.
If for each $x$ in $[-5,5]$, $f(-x)=-f(x)$, i.e., $f(x)$ is an
odd function, then ${\displaystyle \int_{-5}^{5}f(x)\,dx}=$

\begin{multicols}{4}
\begin{enumerate}
\item 26\columnbreak
\item 13\columnbreak
\item 0\columnbreak
\item -13\columnbreak
\end{enumerate}
\end{multicols}\vfill\dspace
\item The average value of a function $f$ on the interval $[a,b]$ is given
by 

\begin{multicols}{4}
\begin{enumerate}
\item $\frac{1}{b-a}{\displaystyle \int_{a}^{b}f(x)\,dx}$\columnbreak
\item $(b-a){\displaystyle \int_{a}^{b}f(x)\,dx}$\columnbreak
\item $\frac{b}{a}\int_{a}^{b}f(x)\,dx$\columnbreak
\item $\frac{f(b)-f(a)}{b-a}$\columnbreak
\end{enumerate}
\end{multicols}\vfill\dspace\newpage
\item If ${\displaystyle \int_{1}^{2}f(x)\,dx}=5$, then ${\displaystyle \int_{1}^{2}\left(f(x)+2x\right)\,dx}$
is equal to

\begin{multicols}{4}
\begin{enumerate}
\item $5$\columnbreak
\item $8$\columnbreak
\item $10$\columnbreak
\item $2x\cdot f(x)$\columnbreak
\end{enumerate}
\end{multicols}\vspace{2.25cm}
\item If ${\displaystyle \int_{b}^{a}f(x)\,dx}=13$, then ${\displaystyle \int_{a}^{b}f(x)\,dx}=$

\begin{multicols}{4}
\begin{enumerate}
\item $13$\columnbreak
\item $6.5$\columnbreak
\item $-13$\columnbreak
\item $0$\columnbreak
\end{enumerate}
\end{multicols}\dspace
\item The volume of the solid generated by revolving the shaded region $R$
about the $x$-axis is computed using the formula 

\begin{multicols}{3}
\begin{enumerate}
\item $\pi{\displaystyle \int_{1}^{3}\left(f(x)\right)^{2}\,dx}$\vfill
\item $\pi{\displaystyle \int_{1}^{3}f(x)\,dx}$\vfill\columnbreak
\item $\int f(x)\,dx$\vfill
\item $\pi{\displaystyle \int_{1}^{3}\sqrt{1+\left(\frac{dy}{dx}\right)^{2}}\,dx}$\vfill\columnbreak
\item []\renewcommand{\Fx}{x^(1/2)}
\renewcommand{\xmax}{5}
\renewcommand{\xmin}{-1}
\renewcommand{\xstep}{0} %must be xmin+n
\renewcommand{\ymax}{2}
\renewcommand{\ymin}{-.35}
\renewcommand{\ystep}{-} %must be ymin+n
%%%%%%%
\begin{tikzpicture}
\begin{axis}[height=4cm, 
	width=5cm, 
	axis lines=middle,
	%axis line style={-latex}, 
	xmin=\xmin,xmax=\xmax, 
	ymin=\ymin,ymax=\ymax,
	%restrict y to domain=-1:10,
	no markers, 
	ticks=none,
	xlabel={$x$},ylabel={$y$}, 
	%every axis x label/.append style={anchor=north}, 
	%every axis y label/.append style={anchor=east}, 
	area style,tension=0.08]

	\addplot+[domain=1:3,fill,black!30!white]{\Fx} \closedcycle;
	\addplot+[domain=1:3,fill,black!30!white]{\Fx} \closedcycle;

	\draw[thick,black,dashed] (axis cs:1,0) node[below] {$1$} -- (axis cs:1,1) ;
	\draw[thick,black,dashed] (axis cs:3,0) node[below] {$3$} -- (axis cs:3,1.73205);
	
	\addplot[domain=0:\xmax,thick,smooth,red,samples=\samples]{\Fx} node[above,pos=.35,yshift=3pt] {$f(x)$}; 
	
	\node at (axis cs:2,.75) {$R$};
\end{axis}
\end{tikzpicture}
\end{enumerate}
\end{multicols}\dspace
\item ${\displaystyle \int_{3}^{3}}\frac{x^{4}+x^{2}+1}{\sqrt{x^{8}+x^{4}+2}}$

\begin{multicols}{4}
\begin{enumerate}
\item 9\columnbreak
\item 0\columnbreak
\item 6\columnbreak
\item $\frac{1}{\sqrt{3}}$\columnbreak
\end{enumerate}
\end{multicols}\dspace
\end{enumerate}
Find each indefinite integral for problems \ref{enu:first-indefinite-int-prob}
through \ref{last-indef-int-prob}. \emph{Show your work.}
\begin{enumerate}[resume]
\item \label{enu:first-indefinite-int-prob}${\displaystyle \int\frac{dx}{(x-1)(x+3)}}$\vfill
\item ${\displaystyle \int xe^{-x}\,dx}$\vfill\newpage
\item ${\displaystyle \int\frac{dx}{9+x^{2}}}$\vfill
\item ${\displaystyle \int\frac{2x}{9+x^{2}}\,dx}$\vfill
\item \label{last-indef-int-prob}${\displaystyle \int\tan^{5}x\sec^{2}\,dx}$\vfill
\item Find the area of the region bounded by the graphs of $y=x^{2}+2$
and $y=x-2$ from $x=-1$ to $x=1$.\vfill\newpage
\item Find the volume of the solid generated by revolving the region bounded
by the the $x$-axis and the curve $y=\sqrt{x+1}$ from $x=0$ to
$x=3$ around the $x$-axis.\vfill
\item Find the length of the arc on the graph of $y=\frac{2}{3}x^{\nicefrac{3}{2}}+9$
from the point $(0,9)$ to the point $(3,2\sqrt{3}+9)$.\vfill
\item If $F(x)$ is an antiderivative of $f(x)$ on the interval $[a,b]$
and $F(a)=2$, $F(b)=5$ then ${\displaystyle \int_{a}^{b}f(x)\,dx}=$

\begin{multicols}{4}
\begin{enumerate}
\item 7\columnbreak
\item 5\columnbreak
\item 2\columnbreak
\item 3\columnbreak
\end{enumerate}
\end{multicols}
\item Solve the initial value problem, $\frac{dy}{dx}=\sqrt{x+1}$ and $y(0)=1$.\vfill
\item Evaluate ${\displaystyle \int_{-1}^{3}\left|x-2\right|}\,dx$.\vfill
\item Use the figure to evaluate ${\displaystyle \int_{-4}^{4}f(x)\,dx}$.

\begin{enumerate}
\item []\hfill\renewcommand{\Fx}{x+4}
\renewcommand{\Gx}{-2*x+4}
\renewcommand{\xmax}{7}
\renewcommand{\xmin}{-7}
\renewcommand{\xstep}{0} %must be xmin+n
\renewcommand{\ymax}{5}
\renewcommand{\ymin}{-5}
\renewcommand{\ystep}{-} %must be ymin+n
%%%%%%%
\begin{tikzpicture}
\begin{axis}[%height=5cm, 
	width=6cm,
	axis equal,
	axis lines=middle,
	%axis line style={-latex}, 
	xmin=\xmin,xmax=\xmax, 
	ymin=\ymin,ymax=\ymax,
	%restrict y to domain=-1:10,
	%no markers, 
	%ticks=none,
	%minor grid style={gray!25},
	%major grid style={black!50},
	minor x tick num={4},
	minor y tick num={4},
	grid=both,
	xlabel={$x$},ylabel={$y$}, 
	%every axis x label/.append style={anchor=north}, 
	%every axis y label/.append style={anchor=east}, 
	]

	%\addplot+[domain=1:3,fill,black!30!white]{\Fx} \closedcycle;
	%\addplot+[domain=1:3,fill,black!30!white]{\Fx} \closedcycle;

	%\draw[thick,black,dashed] (axis cs:1,0) node[below] {$1$} -- (axis cs:1,1) ;
	%\draw[thick,black,dashed] (axis cs:3,0) node[below] {$3$} -- (axis cs:3,1.73205);
	
	\addplot[domain=-4:0,thick,smooth,red,samples=\samples]{\Fx} node[above,pos=.35,yshift=10pt] {$f(x)$};
	\addplot[domain=0:4,thick,smooth,red,samples=\samples]{\Gx} node[above,pos=.35,yshift=5pt] {$$}; 
	
	\node at (axis cs:2,.75) {$R$};
\end{axis}
\end{tikzpicture}
\end{enumerate}
\end{enumerate}

\end{document}
